% Section 10.9, from lectures/L10/README.md.

\section{Summary}
\label{c10:sec:review}

\needspace{6\baselineskip}
\subsection*{What you should be able to do now}
\begin{itemize}
  \item Say why communication protocols use frames, checksums and arbitration.
  \item Describe the CAN frame's format field by field, and construct a complete frame by hand.
  \item Explain how wired-AND arbitration resolves contention without a bus master, and decide which
    of two nodes wins.
  \item Apply the bit stuffing rule to a bit sequence, by hand.
  \item Read \file{can_def.vhd}: say where every constant, subtype and the CRC polynomial come from,
    and why the constants that can be computed are computed rather than typed in.
\end{itemize}

\needspace{6\baselineskip}
\subsection*{Questions to test yourself with}
\begin{enumerate}
  \item Why does a protocol need a length field when the frame has a clear beginning anyway?
  \item Why does CAN have no destination address, and what does it use instead?
  \item Two nodes start transmitting at the same time. Which wins, and what does the loser do?
  \item Why must a transmitter insert a bit after five identical ones in a row, and why does the
    rule not apply to the whole frame?
  \item What happens to the CRC computation for a stuffed bit?
  \item Why is the sample point around 70~\% into the bit period instead of in the middle?
\end{enumerate}

\needspace{6\baselineskip}
\subsection*{Target}
By the end of the session, everyone in the group should be able to say where the constants in
\file{can_def.vhd} come from. It is a short file and the rest of the project reads it, so it is worth
understanding straight away.

The targets in these chapters are just that, targets. They are not milestones, they are not graded,
and the group decides for itself in which order the modules are built; see
appendix~\ref{app:project}.

\needspace{6\baselineskip}
\subsection*{Looking ahead}
\Chapref{c11:ch} goes from protocol to block diagram: the controller's architecture, the top level
every later module is instantiated in, the register map the driver class will consume, and the
project's simulation flow in GHDL.
